\documentclass[UTF8,a4paper,10pt,fontset=none]{ctexart}
\usepackage[left=16.5mm,right=16.5mm,top=17mm,bottom=17mm]{geometry}
% Compatibility with the locally installed pre-2025 LaTeX format.
\providecommand{\IfDocumentMetadataT}[1]{}
\usepackage{amsmath,amssymb,graphicx,booktabs,tabularx,array,xcolor,fancyhdr,hyperref}
\setmainfont{Times New Roman}
\setsansfont{Arial}
\setCJKmainfont[BoldFont=Microsoft YaHei]{SimSun}
\setCJKsansfont{Microsoft YaHei}
\definecolor{accent}{HTML}{176B65}
\definecolor{ink}{HTML}{172C35}
\definecolor{muted}{HTML}{4D5B64}
\hypersetup{colorlinks=true,urlcolor=accent,linkcolor=accent}
\setlength{\parindent}{0pt}
\setlength{\parskip}{5pt}
\renewcommand{\arraystretch}{1.32}
\newcolumntype{Y}{>{\raggedright\arraybackslash}X}
\pagestyle{fancy}\fancyhf{}
\fancyhead[L]{\small\sffamily COT / GHI · 图表精修评审版}
\fancyhead[R]{\small\sffamily 30 OCT 2026}
\fancyfoot[L]{\small\color{muted}真实数据 · 来源明确 · 逐图评阅}
\fancyfoot[R]{\small\thepage}
\renewcommand{\headrulewidth}{0.3pt}
\setlength{\headheight}{14pt}
\setlength{\emergencystretch}{2em}
\begin{document}\color{ink}
{\color{accent}\sffamily\bfseries\fontsize{27}{32}\selectfont COT论文\quad 图表精修第一版}\par
{\sffamily\large 从原文设计逻辑，到可复核的模型与结果表达}\par
\vspace{3mm}
本版先交付可以逐页评阅的结构图、流程图、真实云事件及结果图表。图形来自实际代码和保存预测，没有添加模拟结果。\textbf{它是图表评审版，不是已经完成独立验证的投稿稿。} 湖北同期单帧主结果、同一来源的COT质量及厚度分层仍需补齐；这些位置不借用旧分数填充。

\textbf{表1\quad 从原论文借鉴什么，怎样转化为本研究的图表。}\par
\begin{tabularx}{\linewidth}{p{49mm}YY}\toprule
原文与图表 & 作者为何需要它 & 本版采用的表达\\\midrule
Chen et al., Applied Energy (2023), Fig.4 / Table3 & 解释多通道如何编码、融合；表中保持站点与方法的明确对应 & 结构展开到算子和张量；精确数值与比较问题分表\\
Papachristopoulou et al., AMT (2024), Fig.2 / 11 / 15 & 把物理关系、误差、预报时效及云场变化分开回答 & 绝对误差与配对区间并列；COT输入质量不替代GHI收益\\
Zhang et al., GRL (2026), Fig.1 / 2 / 3 & 连接网络、逐时效表现与空间演变 & 网络和递推分开；少量可读关键帧连接连续轨迹\\
用户提供的Baguan-solar结构图 & 两阶段清楚分工；多源输入与关键融合展开 & 区分数据流程与网络结构；只画本项目实际运算\\\bottomrule\end{tabularx}
原论文截图、页码和逐项分析保留在配套阅读册。这里借鉴论证和布局，不复制其数值，也不将跨论文RMSE当作同一测试集排名。

\vspace{3mm}\textbf{表2\quad 本版图表使用的样本范围。}\par
\begin{tabularx}{\linewidth}{p{43mm}p{30mm}Y}\toprule
内容 & 有效行 / 日期 & 用途与边界\\\midrule
湖北四臂预报 & 15,275 / 22 & 三站、三seed；验证集；历史COT与新增未来COT分别比较\\
AGRI × COT来源替换 & 8,111 / 12 & 三站、三seed；已有共同可替换掩膜；开发期固定头诊断\\
条件信息探针 & 32,838 / 50 & 开发日期折外拟合；三份冻结基线来源，不是三次新探针训练\\
真实云事件 & 2025-08-06 & 日历规则选例；牛首13个有效GHI区间；解释现象而非总体收益\\
湖南单帧补充 & 4,440 / 51 & 两站、三seed；固定权重重评分，沿用原太阳几何处理\\\bottomrule\end{tabularx}
{\small\color{muted}行数不是独立天气样本数。不同表的日期、模型和输入范围不同，不能串成一条“精度逐步改善”的曲线。所有含未来真实图像的结果均明确标为离线诊断。\par}
\clearpage\noindent{\color{accent}\sffamily\bfseries\fontsize{23}{28}\selectfont 01}\quad{\sffamily\bfseries\Large 核心网络结构}\par\vspace{2mm}{\color{muted}建议放在方法节：回答“COT在网络中到底怎样起作用”。}\par\vspace{4mm}
\begin{center}\includegraphics[width=177mm]{images/architecture.pdf}\end{center}
{\small\color{muted}按实际代码展开反演、逐像元编码、空间位置保留、融合及GHI读出。湖南与湖北的宽度和辅助槽分别标明，不能混用。结构已核对；最终湖北同期主结果仍待与权重完整绑定。参数与空间解释见方法附表。\par}
\clearpage\noindent{\color{accent}\sffamily\bfseries\fontsize{23}{28}\selectfont 02}\quad{\sffamily\bfseries\Large 从数据到GHI的完整路径}\par\vspace{2mm}{\color{muted}建议放在方法节开头，与上一页结构图各司其职。}\par\vspace{4mm}
\begin{center}\includegraphics[width=177mm]{images/workflow.pdf}\end{center}
{\small\color{muted}同期估计与可部署预报分开；预报只能从初始化之前的观测出发。CPP只作监督和参考评价。区域图、反演上下文、站点邻域尺度不同；GHI对应以目标时刻结束的15分钟平均。未来真实AGRI仅用于另列的离线诊断。\par}
\clearpage\noindent{\sffamily\bfseries\Large 方法附表：实现差异与共同控制}\par\vspace{4mm}
\textbf{表3\quad 代码实测的规模差异。}\par
\begin{tabularx}{\linewidth}{p{46mm}Yr}\toprule
实现 & 结构区别 & 参数数\\\midrule
湖南COT反演器 & $w=32$；$32\to64\to128$ & 470,913\\
湖北COT反演器 & $w=16$；$16\to32\to64$ & 119,233\\
湖南AGRI基线 / 加COT & 读出输入$288/324$；COT为$1\to4$带偏置 & 22,177 / 24,489\\
湖南辅助 / 反演特征 & $32\to4$；两者相同容量 & 24,613\\
湖北四臂单帧设计 & COT与云概率双槽$2\to4$，不含偏置 & 24,489\\\bottomrule\end{tabularx}
{\small\color{muted}参数计数只描述对应网络；不能把两个模型参数数相近等同于整个系统训练资源匹配。\par}
\textbf{时间与输入边界。} 同期估计可用目标时刻真实AGRI；预报先从已知历史预测未来AGRI，再经冻结反演器得到未来COT。GHI标签是以目标时刻结束的15分钟平均；像元时刻配准不等于瞬时像元与地面平均实现严格辐射闭合。

\textbf{不能省略的对照差异。} 湖南的AGRI辅助特征与冻结反演器特征控制彼此等容量，但比显式COT臂多124个可训练参数。湖北四臂使用相同辅助槽和参数形状，通过中性零槽定义对照。两种设计都需要保留各自身份，不能合并成“全部模型严格等容量”。

\textbf{空间结构怎样进入网络。} 单帧头的$1\times1$卷积逐像元编码，不在这一步混合相邻位置。随后拼接并保留通道、行、列顺序展平，才由全连接读出联合九个位置。因此它保留空间位置，却没有显式建模运动或历史轨迹；不能仅因输入整张小图就写成已学到云输运。

\textbf{反演上下文与GHI邻域。} R在$16\times16$范围获取上下文，GHI读取同站点$3\times3$局地特征。加入COT的收益可能同时受云产品监督和更大上下文影响。要主张光学厚度本身的独特贡献，还应有相同视场与训练资源的表示对照。

\clearpage\noindent{\color{accent}\sffamily\bfseries\fontsize{23}{28}\selectfont 03}\quad{\sffamily\bfseries\Large 历史信息和未来信息分别有没有用}\par\vspace{2mm}{\color{muted}建议放预报应用节：先给真实GHI，再解释COT。}\par\vspace{4mm}
\begin{center}\includegraphics[width=177mm]{images/forecast.pdf}\end{center}
\textbf{表4\quad 湖北三站逐站GHI RMSE（W m$^{-2}$）。}\par
\begin{tabularx}{\linewidth}{Yrrrr}\toprule
站点 & 无COT & 历史COT & 历史+控制修正 & 历史+预测COT\\\midrule
牛首 & 97.017 & 96.299 & 96.665 & 96.701\\
王寨 & 106.183 & 102.892 & 103.165 & 103.215\\
米庄 & 113.586 & 111.819 & 111.519 & 111.510\\
站点等权 & 105.595 & 103.670 & 103.783 & 103.808\\\bottomrule\end{tabularx}
{\small\color{muted}三seed先分别计算RMSE，再按站点和seed等权；不是把三seed预测集成后评分。后两臂具有相同参数量，并共用已训练的历史COT基底。\par}
\textbf{关键判断。} 历史COT点差为$-1.925$ W m$^{-2}$，日期块95\%区间$[-4.734,+0.670]$；新增未来COT相对同容量控制为$+0.026$，区间$[-0.003,+0.059]$。两者均没有给出可靠的正向增益。

\clearpage\noindent{\color{accent}\sffamily\bfseries\fontsize{23}{28}\selectfont 04}\quad{\sffamily\bfseries\Large 真实云图怎样连接COT与辐照度}\par\vspace{2mm}{\color{muted}机制案例；不能将个例改善替代总体预报收益。}\par\vspace{4mm}
\begin{center}\includegraphics[width=174mm]{images/event.pdf}\end{center}
{\small\color{muted}2025年8月6日牛首，初始化12:00 BJT。真实与预测AGRI各紧邻其反演COT；统一C03/C02/C01、反射率0–1平方根显示，保留16×16原像元，标出站点与3×3输入。下方COT和15分钟GHI共享时间轴；缺测断线。灰线已含历史COT，紫线用真实未来图COT，仅作离线诊断。三份冻结基线来源的曲线取展示均值；R-COT超过100仍保留。本例按日历有效性选择，非按收益挑选。\par}
\clearpage\noindent{\color{accent}\sffamily\bfseries\fontsize{23}{28}\selectfont 05}\quad{\sffamily\bfseries\Large 区分底图误差与COT映射的响应}\par\vspace{2mm}{\color{muted}同一固定模型的输入替换；独立于上一页个例和正式验证主表。}\par\vspace{4mm}
\begin{center}\includegraphics[width=177mm]{images/replacement.pdf}\end{center}
\textbf{表5\quad 固定头的MSE敏感性分解（W$^2$ m$^{-4}$）。}\par
\begin{tabularx}{\linewidth}{Yrr}\toprule
项 & 点估计 & 日期块95\%区间\\\midrule
$E_{FO}-E_{OO}$：图像路径 & +2778.1 & [+938.0, +5228.9]\\
$E_{OF}-E_{OO}$：COT路径 & -5.4 & [-48.9, +33.5]\\
$E_{FF}-E_{FO}-E_{OF}+E_{OO}$：交互 & -51.6 & [-109.2, -7.8]\\
$E_{FF}-E_{OO}$：总变化 & +2721.1 & [+920.4, +5105.2]\\\bottomrule\end{tabularx}
{\small\color{muted}O为真实图来源、F为预测图来源。前两项与交互之和等于总变化，区间本身不能相加。这里是固定映射对输入替换的响应，不是独立物理误差预算。\par}
\textbf{可成立的解释。} 换成真实AGRI明显降低误差，但在同源真实输入下额外COT仍未改善这个固定头。因此底图信息损失是重要线索，不能据此宣称唯一原因，也不能把映射适应问题排除。

\clearpage\noindent{\color{accent}\sffamily\bfseries\fontsize{23}{28}\selectfont 06}\quad{\sffamily\bfseries\Large 真实图COT有信息，预测源保留了多少}\par\vspace{2mm}{\color{muted}条件探针使用开发日期折外拟合，不能与正式验证绝对RMSE混排。}\par\vspace{4mm}
\begin{center}\includegraphics[width=177mm]{images/probe.pdf}\end{center}
\textbf{先看参照。} 已有历史COT预报基线为118.799 W m$^{-2}$；上下文探针为124.850，加入预测COT为124.307，两者都比已有基线差。只有使用真实未来图COT的离线探针达到99.569。

\textbf{再看来源替换。} 对已经拟合的探针更换COT来源，可检验其响应是否随来源变化。图中的虚线列依赖未来观测。较好的离线分数说明这种来源在当前诊断中可被利用，却不证明已经有可部署的改善方案。

{\small\color{muted}32,838条共同记录、50个开发目标日期；三个站点与三份冻结基线来源等权。探针为确定性岭回归，不能把这三份来源写成三次独立探针种子。误差线为5,000次目标日期块配对bootstrap的95\%区间。底部纵轴是GHI误差差，不是COT物理精度。\par}
\clearpage\noindent{\color{accent}\sffamily\bfseries\fontsize{23}{28}\selectfont 07}\quad{\sffamily\bfseries\Large 区域差异与同期证据的适用范围}\par\vspace{2mm}{\color{muted}两张材料各说明一个问题：同日三站差异，以及旧处理下的固定权重结果。}\par\vspace{4mm}
\begin{center}\includegraphics[width=170mm]{images/stations.pdf}\end{center}
{\small\color{muted}同一事件三站分别有13、14、11个有效GHI区间；各排共用尺度，缺测不连线。真实未来COT仍仅作离线诊断，不把三站当作独立新天气事件。\par}
\vspace{3mm}\textbf{表6\quad 湖南单帧固定权重重评分（补充诊断；W m$^{-2}$）。}\par
\begin{tabularx}{\linewidth}{Yrrrr}\toprule
站点 & AGRI & AGRI+COT & 辅助特征 & 反演特征\\\midrule
四里 & 103.558 & 99.877 & 103.457 & 99.962\\
竺家 & 112.805 & 108.113 & 112.571 & 106.950\\
站点等权 & 108.181 & 103.995 & 108.014 & 103.456\\\bottomrule\end{tabularx}
\textbf{如何读。} 两站总体108.181$\to$103.995，点降幅3.870\%；反演隐特征总体103.456，说明这组结果不能证明显式COT是唯一有效表示。

{\small\color{muted}此表保留原权重和原太阳几何处理；剔除已确认无效站点日后，在4,440行、51日期上统一重评分。模型曾按剔除前验证集选模，不是重训、独立测试或修正太阳几何后结果，不能填充湖北同期主结果缺口。\par}
\clearpage\noindent{\color{accent}\sffamily\bfseries\fontsize{23}{28}\selectfont 08}\quad{\sffamily\bfseries\Large 进入正文前，还要补齐什么}\par\vspace{2mm}{\color{muted}保留科学问题和图表位置；不以美化代替缺失证据。}\par\vspace{4mm}
\textbf{表7\quad 图表与论证的对应关系。}\par
\begin{tabularx}{\linewidth}{p{33mm}YY}\toprule
问题 & 本版已完成 & 仍需完成\\\midrule
方法如何复现 & 数据流程、实际网络层次、空间尺度、冻结与监督边界 & 将最终同期主结果的权重、配置与图中的区域实现一一绑定\\
COT能否改善同期GHI & 湖南固定权重补充表；可复核的代码与方法图 & 湖北同一修正来源的同期主比较、容量/表示控制\\
哪些云况受益 & 真实云事件与三站响应展示 & 同一最终来源的CPP质量、厚度分层及另一互补事件\\
未来COT是否有效 & 湖北四臂、逐时效、日期块区间 & 这些是验证证据；独立确认须另按冻结方案进行\\
为什么未来收益有限 & 固定头替换与折外探针分开展示 & 不能给唯一归因或将不同队列的差值画成收益衰减量\\\bottomrule\end{tabularx}
\textbf{建议阅读顺序。} 主文先用完整路径和网络结构解释方法；同期主结果和厚度分层补齐后，连接真实云事件；最后用预报表与逐时效图检验应用边界。固定头来源替换可放预报节或讨论；条件探针和三站完整曲线适合补充材料。

\textbf{统一视觉规则。} 观测用黑色；已有基线用灰色；历史COT用青绿色；预测COT用橙色；未来真实来源用紫色并标离线。差值始终为候选减参照，负值改善。空间图保留真实像元，轨迹有缺测就断线，不为平滑和锐化改变数据。

\vspace{3mm}\textbf{原文与工具来源}\par
{\small
Chen et al. (2023), \textit{Applied Energy}, 352, 121979. \href{https://doi.org/10.1016/j.apenergy.2023.121979}{DOI}\par
Papachristopoulou et al. (2024), \textit{AMT}, 17, 1851--1877. \href{https://amt.copernicus.org/articles/17/1851/2024/}{全文}\par
Zhang et al. (2026), \textit{GRL}, 53, e2025GL119553. \href{https://doi.org/10.1029/2025GL119553}{DOI}\par
Baguan-solar, 用户结构截图参考。\href{https://arxiv.org/abs/2603.14845}{arXiv:2603.14845}\par
科学可视化技能用于设计与审图；数值由保存预测重算。Kassis et al. (2026), \textit{Scientific Agent Skills: A Library of Procedural Knowledge for Research Agents}. \href{https://arxiv.org/abs/2609.00065}{arXiv:2609.00065}。技能不是科学结论的证据。\par}
\end{document}
